\documentclass[12pt,a4paper]{article}

\usepackage[T1]{fontenc}
\usepackage[utf8]{inputenc}
\usepackage[english]{babel}
\usepackage{lmodern}
\usepackage{amsmath,amssymb}
\usepackage[hidelinks]{hyperref}
\usepackage{graphicx}

\makeatletter
\setlength{\@fptop}{0pt}
\setlength{\@fpbot}{0pt plus 1fil}
\makeatother

\title{Fairness Evaluation of Classification Models}
\author{}
\date{}

\begin{document}
\maketitle

\section{Fairness as a Measure of Trustworthiness}

\subsection{Foundational Concepts}

An aggregate performance score can hide differences in how a classifier treats groups of examples. Group fairness evaluation makes these differences visible by comparing prediction outcomes or error rates across a specified attribute. There is no single definition that captures every relevant notion of fairness; different definitions answer different questions and can lead to different conclusions \cite{verma2018}. Let $Y$ denote the true outcome, $\widehat{Y}$ the predicted outcome, and $A$ a group attribute. For binary classification, the true-positive and false-positive rates in group $a$ are
\begin{equation}
    \operatorname{TPR}_a=\Pr(\widehat{Y}=1\mid Y=1,A=a),
    \qquad
    \operatorname{FPR}_a=\Pr(\widehat{Y}=1\mid Y=0,A=a).
    \label{eq:group-error-rates}
\end{equation}

Hardt, Price, and Srebro define \emph{equalized odds} by requiring both rates to be equal across groups; \emph{equal opportunity} requires only equal TPR \cite{hardt2016}:
\begin{align}
    \operatorname{TPR}_a&=\operatorname{TPR}_{a'}
    \quad\text{and}\quad
    \operatorname{FPR}_a=\operatorname{FPR}_{a'}
    &&\text{(equalized odds)}, \label{eq:equalized-odds}\\
    \operatorname{TPR}_a&=\operatorname{TPR}_{a'}
    &&\text{(equal opportunity)}. \label{eq:equal-opportunity}
\end{align}
These are group-level statistical criteria: they compare observed labels, predictions, and group membership, but do not explain why a difference arose or establish that a use is socially acceptable. Excluding $A$ from a model's inputs is not sufficient to ensure similar outcomes, since other inputs may carry information about group membership \cite{hardt2016}. The criteria are therefore measurements to interpret in context, rather than a universal verdict about a model.

\subsection{Complementary Metrics}

Equalized odds and equal opportunity condition on the true outcome. Other measures compare different quantities and should be reported separately. The \emph{selection rate} for group $a$ is the proportion receiving a positive prediction,
\begin{equation}
    \operatorname{SR}_a=\Pr(\widehat{Y}=1\mid A=a).
    \label{eq:selection-rate}
\end{equation}
The statistical parity difference compares selection rates between groups, while a disparate impact ratio compares a group's rate with a caller-selected reference group's rate. These measures do not condition on $Y$: a difference may reflect outcome prevalence as well as model behavior. They consequently provide a different view from error-rate comparisons \cite{hardt2016,verma2018}.

Predictive parity examines the correctness of positive predictions through positive predictive value (PPV), $\Pr(Y=1\mid\widehat{Y}=1,A=a)$. Negative predictive value (NPV), $\Pr(Y=0\mid\widehat{Y}=0,A=a)$, gives the corresponding view for negative predictions. Comparing these with TPR, FPR, and selection rate makes clear which kinds of outcomes differ. The metric should be tied to the task's positive class: changing that class changes the interpretation of all class-conditional rates.

These criteria are not interchangeable, and some cannot generally be satisfied together. Chouldechova and Kleinberg, Mullainathan, and Raghavan show incompatibilities among calibration and error-balance conditions under specified assumptions, including settings with differing group base rates and imperfect prediction \cite{chouldechova2017,kleinberg2017}. Pleiss et al. likewise analyze the tension between calibrated risk estimates and balanced error rates \cite{pleiss2017}. These results do not imply that any one criterion is always preferable; they motivate presenting multiple measurements with their definitions and task assumptions rather than reducing them to one score.

\subsection{Scores, ROC Curves, and Calibration}

For a model that returns scores, threshold-dependent TPR and FPR can be plotted as a receiver operating characteristic (ROC) curve. An overall ROC curve describes discrimination between positive and negative examples. Group-conditional ROC curves instead compare this discrimination separately by group. Hardt et al. relate the latter curves to equalized odds across thresholds \cite{hardt2016}. The area between group ROC curves (ABROCA) summarizes their difference over the ROC space; it complements per-group area under the curve (AUC), since two groups can have similar AUC values but differently shaped curves \cite{gardner2019}. Neither score-based summary alone establishes parity at a particular operating threshold.

Probability quality can also vary across groups. Brier score and log loss measure probabilistic prediction error and can be compared across groups as complementary performance measures. They are not fairness criteria by themselves. This work evaluates calibration separately with adaptive calibration error (ACE); ECE is therefore not duplicated in the fairness report. Calibration and error-rate parity remain distinct properties, and their trade-offs depend on the setting \cite{pleiss2017}.

\subsection{Intersectional Evaluation and Reporting}

An audit limited to one attribute at a time can conceal differences affecting intersections of attributes. In a study of commercial gender-classification systems, Buolamwini and Gebru reported substantially different error rates across groups formed by gender and skin type \cite{buolamwini2018}. This motivates evaluating intersectional categories where relevant and supported by the data. Such finer slices can contain few examples, so each reported rate should be accompanied by its group count and interpreted with care.

Model Cards propose documenting a model's intended use, evaluation conditions, and performance across relevant groups \cite{mitchell2019}. In this work, fairness measurements are reported as one part of a broader trustworthiness profile alongside predictive performance, calibration, and other evaluations. The report identifies which attributes and comparisons were supplied, which metrics were calculated, and which dimensions were not measured. It does not classify a model as fair or unfair.

\section{Fairness Evaluation in This Project}

The evaluation function operates on labels, predictions, and group attributes supplied in sample order. It does not depend on a particular model architecture or dataset, so the same classification metrics can be used with image, text, or tabular models if those aligned values are available. Group labels are not inferred by the module. The choice and quality of attributes, the definition of the target, and the intended use remain part of the evaluation design.

For each supplied categorical attribute, the report records group counts, missing attribute values, and the true and predicted class distributions. It calculates accuracy, balanced accuracy, macro precision, macro recall, and macro F1 for each group. For each class it reports selection rate, TPR, FPR, false-negative rate, true-negative rate, PPV, NPV, and F1. The across-group summaries include maximum-minus-minimum ranges for these values. If the caller supplies a reference group, the report also gives signed differences for statistical parity, equal opportunity, FPR, predictive parity, and negative predictive parity, as well as a selection-rate ratio and average odds difference. The reference group is explicit; no preferred group is chosen by the implementation.

When class probabilities are provided, the report adds group-level Brier score and log loss, one-versus-rest AUC, and pairwise ABROCA for groups with both positive and negative examples for the class. These quantities describe probabilistic prediction error and score separation; they are kept distinct from threshold-based gaps. Calibration error is measured separately with ACE in the calibration module. If an accepted-sample mask is supplied, the report includes group coverage and accuracy among accepted samples.

All results are descriptive estimates on the evaluated data. Missing group values are excluded from that attribute's group comparisons, and a rate may be undefined if its denominator is zero. The implementation warns about small groups but does not calculate confidence intervals or significance tests, so observed gaps may be imprecise. It also does not measure sensitive attributes that were not supplied, causal explanations, individual or counterfactual fairness, or fairness for regression and ranking tasks. These limitations and the selected evaluation population should accompany interpretation of the metrics.

\subsection{CIFAR-10 and the Availability of Group Data}

CIFAR-10 contains object-class labels but no demographic or other independent group attributes. An earlier experiment hard-coded the groups \emph{animals} and \emph{vehicles} from the ground-truth class. That partition is not a valid fairness comparison: for several one-versus-rest outcomes, one group contains no positive examples, so group TPR and ROC comparisons are undefined. Its accuracy difference of $0.0617$ describes a class-composition profile only and must not be interpreted as a fairness result.

The CIFAR-10 runner now requires external group metadata aligned row-for-row with the test examples. More generally, the shared audit function accepts predictions and group attributes from any classification model or dataset. If suitable group data are unavailable, the evaluation can report predictive performance, but it cannot make a group fairness comparison.

\section{Implementation of the Fairness Module}

The metric calculation is implemented in \path{evaluation/fairness.py} and exposed through \path{experiments/fairness.py}. Its public function takes ground-truth labels, predictions, one or more group attributes, and optional probabilities. It returns a structured report that preserves the group and class measurements. The report contains no aggregate fairness score, ranking, threshold, or fair/unfair verdict.

The reusable metric calculation accepts labels, predictions, categorical group attributes, and optional class probabilities. It does not load a model or dataset. The model-agnostic CSV adapter in \path{experiments/fairness_csv.py} reads these aligned outputs from any classifier. The CIFAR-10 runner in \path{experiments/fairness_cifar10.py} is a model-and-dataset-specific adapter; it requires a CSV containing real group attributes aligned with the test split and does not derive groups from target labels.

To run the generic CSV adapter from the project root, provide columns for labels, predictions, group attributes, and optionally probabilities:
\begin{verbatim}
.\.venv\Scripts\python.exe -m experiments.fairness_csv `
    --input predictions.csv --label-column y_true --prediction-column y_pred `
    --group-columns sex age_band --output results/fairness.json
\end{verbatim}

For a useful audit, select measures that reflect the intended task, preserve sample alignment between predictions and group attributes, and report group sizes and outcome distributions alongside the rates. The purpose of this module is to make observed differences and unmeasured dimensions visible so that model suitability can be considered for a specific use.

\begin{thebibliography}{99}

\bibitem{hardt2016}
M. Hardt, E. Price, and N. Srebro.
\newblock Equality of Opportunity in Supervised Learning.
\newblock In \emph{Advances in Neural Information Processing Systems 29}, pp. 3315--3323, 2016.
\newblock \url{https://papers.neurips.cc/paper_files/paper/2016/file/6a9659feb1216f14f7384ba499518b38-Paper.pdf}.

\bibitem{verma2018}
S. Verma and J. Rubin.
\newblock Fairness Definitions Explained.
\newblock In \emph{Proceedings of the 2018 IEEE/ACM International Workshop on Software Fairness}, pp. 1--7, 2018.
\newblock \url{https://doi.org/10.1145/3194770.3194776}.

\bibitem{chouldechova2017}
A. Chouldechova.
\newblock Fair Prediction with Disparate Impact: A Study of Bias in Recidivism Prediction Instruments.
\newblock \emph{Big Data}, 5(2):153--163, 2017.
\newblock \url{https://doi.org/10.1089/big.2016.0047}.

\bibitem{kleinberg2017}
J. Kleinberg, S. Mullainathan, and M. Raghavan.
\newblock Inherent Trade-Offs in the Fair Determination of Risk Scores.
\newblock In \emph{Proceedings of the 8th Innovations in Theoretical Computer Science Conference}, LIPIcs 67, Article 43, 2017.
\newblock \url{https://doi.org/10.4230/LIPIcs.ITCS.2017.43}.

\bibitem{pleiss2017}
G. Pleiss, M. Raghavan, F. Wu, J. Kleinberg, and K. Q. Weinberger.
\newblock On Fairness and Calibration.
\newblock In \emph{Advances in Neural Information Processing Systems 30}, 2017.
\newblock \url{https://papers.nips.cc/paper/2017/file/b8b9c74ac526fffbeb2d39ab038d1cd7-Paper.pdf}.

\bibitem{gardner2019}
J. Gardner, C. Brooks, and R. Baker.
\newblock Evaluating the Fairness of Predictive Student Models Through Slicing Analysis.
\newblock In \emph{Proceedings of the 9th International Conference on Learning Analytics and Knowledge}, 2019.
\newblock \url{https://doi.org/10.1145/3303772.3303791}.

\bibitem{buolamwini2018}
J. Buolamwini and T. Gebru.
\newblock Gender Shades: Intersectional Accuracy Disparities in Commercial Gender Classification.
\newblock In \emph{Proceedings of the 1st Conference on Fairness, Accountability and Transparency}, PMLR 81, pp. 77--91, 2018.
\newblock \url{https://proceedings.mlr.press/v81/buolamwini18a.html}.

\bibitem{mitchell2019}
M. Mitchell, S. Wu, A. Zaldivar, P. Barnes, L. Vasserman, B. Hutchinson, E. Spitzer, I. D. Raji, and T. Gebru.
\newblock Model Cards for Model Reporting.
\newblock In \emph{Proceedings of the Conference on Fairness, Accountability, and Transparency}, pp. 220--229, 2019.
\newblock \url{https://doi.org/10.1145/3287560.3287596}.

\end{thebibliography}

\end{document}
